% Target: rmp-workbench-iii-diagram-three
% Formula: The north-west and south-east rotated boundary actions agree.
% Ink: Canonical public rotated tenkzfree closure with a lambda insertion.
% Boundary: Each side has four directed external legs and one closed inner ring.
% Source: author source ImagesReview Section III/ImagesReview Section III.tex
%   lines 170-197; archival output Dia3.pdf
% Capabilities: free-graph, rotated-action, ring-closure
\[
  \begin{tenkzfree}\tngroup[frame={rotate=90}]{
    \tnput[circle]{l}{(-6mm,6mm)}{\lambda}
    \tnput[box]{n}{(0mm,10mm)}{} \tnput[box]{w}{(-10mm,0mm)}{}
    \tnput[box]{s}{(0mm,-10mm)}{} \tnput[box]{e}{(10mm,0mm)}{}
    \tnjoin[dir]{n.north}{0mm,18mm} \tnjoin[dir]{w.west}{-18mm,0mm}
    \tnjoin[dir]{0mm,-18mm}{s.south} \tnjoin[dir]{18mm,0mm}{e.east}
    \tnjoin[route=arc, out=0, in=90]{n.east}{e.north}
    \tnjoin[route=arc, out=-90, in=0]{e.south}{s.east}
    \tnjoin[route=arc, out=180, in=-90]{s.west}{w.south}
    \tnjoin[route=arc, out=90, in=180]{w.north}{l.west}
    \tnjoin{l.east}{n.west}
  }\end{tenkzfree}
  =
  \begin{tenkzfree}\tngroup[frame={rotate=90}]{
    \tnput[circle]{l}{(6mm,-6mm)}{\lambda}
    \tnput[box]{n}{(0mm,10mm)}{} \tnput[box]{w}{(-10mm,0mm)}{}
    \tnput[box]{s}{(0mm,-10mm)}{} \tnput[box]{e}{(10mm,0mm)}{}
    \tnjoin[dir]{n.north}{0mm,18mm} \tnjoin[dir]{w.west}{-18mm,0mm}
    \tnjoin[dir]{0mm,-18mm}{s.south} \tnjoin[dir]{18mm,0mm}{e.east}
    \tnjoin[route=arc, out=0, in=90]{n.east}{e.north}
    \tnjoin[route=arc, out=-90, in=0]{e.south}{l.east}
    \tnjoin{l.west}{s.east}
    \tnjoin[route=arc, out=180, in=-90]{s.west}{w.south}
    \tnjoin[route=arc, out=90, in=180]{w.north}{n.west}
  }\end{tenkzfree}.
\]
